\chapter{Methodology} \label{metodologia}
\vspace{-1.9cm}

This chapter details the methodology employed to develop, implement, and evaluate the smart task scheduling strategy for heterogeneous Autonomous Driving Vehicles (ADV) processing systems. To provide a structured and comprehensive overview of the research design, the methodology is divided into the following main stages:

\begin{itemize}
    \item \textbf{Material:} Description of the software, frameworks, and hardware infrastructure utilized to build the simulation and scheduling environment.
    \item \textbf{Setup of the Test Environment:} The integration of the ADV simulator with the task scheduling framework.
    \item \textbf{Design of the ADV Smart Task Scheduling Strategy:} The modeling of the ADV computational workload and the proposed scheduling approach.
    \item \textbf{Implementation and Testing:} The practical application of the proposed strategy within the simulated environment.
    \item \textbf{Results Analysis and Evaluation:} The metrics and benchmarks used to assess the performance of the proposed scheduling algorithm.
\end{itemize}

\section{Material} \label{sec:material}

To develop and evaluate the novel task scheduling strategy, a specific set of tools, frameworks, and hardware components was selected to accurately simulate the computational demands of an Autonomous Driving Vehicle (ADV). This section describes the primary components that comprise the testbed infrastructure.

\subsection{Simulation Platform}
The CARLA (Car Learning to Act) simulator \cite{Dosovitskiy17} was employed as the primary open-source platform to replicate real-world autonomous driving scenarios. CARLA was selected due to its high-fidelity rendering, realistic urban environments, and robust Python API, which provides extensive flexibility in configuring simulation parameters, controlling vehicle sensors (e.g., cameras, LIDAR, radar), and defining complex traffic scenarios involving varying densities, pedestrian behaviors, and weather conditions. Crucially, CARLA allows for the extraction of sensor data in real-time, functioning as the primary workload generator for the ADV perception pipeline.

\subsection{Task Scheduling Framework}
To manage the computational workload generated by the ADV simulator across a heterogeneous processing architecture, the StarPU runtime system \cite{Augonnet2011} was utilized. StarPU is a task programming library for hybrid architectures that provides a unified execution model. It was chosen for its capability to seamlessly manage task graphs (Directed Acyclic Graphs - DAGs) across different processing units (CPUs and GPUs), its comprehensive set of profiling and tracing tools, and, most importantly for this research, its flexible API that allows users to write and plug in custom scheduling policies \cite{starpuscheduling2024}. 

\subsection{Programming Languages and Libraries}
The core integration between the CARLA simulator and the StarPU runtime was developed using \textbf{C++} and \textbf{Python}. Python was heavily utilized for interacting with the CARLA API and managing the high-level simulation scripts, leveraging its extensive ecosystem for data analysis and machine learning. C++ was employed for performance-critical components, specifically for developing the StarPU application client, defining the StarPU \textit{codelets} (task implementations), and interfacing directly with the StarPU C API to ensure minimal overhead during task submission and scheduling.

For the object detection workload within the ADV perception pipeline, the YOLO (You Only Look Once) architecture was incorporated, utilizing deep learning frameworks to process the camera data streams extracted from CARLA. 

\subsection{Test Environment Hardware}
The experiments and simulations were conducted on a heterogeneous computing platform to evaluate the scheduling strategies under realistic hardware constraints. The test environment was composed of a system equipped with both multi-core Central Processing Units (CPUs) and Graphics Processing Units (GPUs), representing the heterogeneous nature of modern ADV embedded systems. [NOTE TO USER: Please insert the specific hardware details here: e.g., Intel Core i9 processor, 32GB RAM, NVIDIA RTX 3080 GPU, running Ubuntu 22.04 LTS].
